\documentclass[UTF8,a4paper]{ctexart}

\usepackage[a4paper,margin=2.2cm]{geometry}
\usepackage{amsmath,amssymb}
\usepackage{booktabs}
\usepackage{tabularx}
\usepackage{enumitem}
\usepackage{xurl}
\usepackage{hyperref}
\usepackage{array}
\usepackage{longtable}
\usepackage{multirow}

\hypersetup{
  colorlinks=true,
  linkcolor=blue,
  urlcolor=blue,
  citecolor=blue
}

\setlist[itemize]{itemsep=3pt,topsep=4pt,leftmargin=2em}
\setlist[enumerate]{itemsep=3pt,topsep=4pt,leftmargin=2em}

\newcommand{\code}[1]{\texttt{#1}}
\newcommand{\gri}{\textbf{GRI}}
\newcommand{\grfm}{\textbf{GRFM}}
\newcommand{\gbr}{\textbf{GBR-AD}}
\newcommand{\gbrvtwo}{\textbf{GBR-AD v2}}

\title{BMAD Essay Research Guideline\quad Revised for GBR-AD v2}
\author{杜博源}
\date{\today}

\begin{document}
\maketitle

\section{实验目标与设置说明}
本文面向医学图像异常检测与定位任务，采用 BMAD 统一框架进行训练与评估。
当前已完成 RESC 数据集上 RD4AD 的第一组对照实验，并已完成 GBR-AD 的 B0 插件化接入与最小闭环验收。
本版 guideline 的修改原则如下。
A 组与 B0 组内容保持不变并继续有效。
从 B1 开始，实验主线升级为 \gbrvtwo，目标是在保留粒计算思想为方法灵魂的前提下，缩小当前版本与顶级会议论文要求之间的差距。

现阶段目标分四层。
第一，形成覆盖 BMAD 六数据集的主结果对比，保证结论具备跨域代表性。
第二，将粒球机制从“可运行模块”升级为“主表示单元 + 主检测单元 + 主证据单元”，构成 \gbrvtwo 的正式方法主线。
第三，用困难子集、消融、效率和证据链分析支撑“为什么有效”而不是只给分数。
第四，形成可直接进入论文正文和附录的图表、公式、可视化与归档材料。

实验原则如下。
训练阶段只使用正常样本。
验证与测试阶段同时包含正常与异常样本。
所有方法使用一致的输入尺寸、预处理流程与评测脚本。
关键结果使用多个随机种子重复运行，并报告均值与标准差。
所有 run 需要保存可复现快照与核心产物索引。
正式主表只接收统一 stopping policy、统一训练预算、统一评测脚本下得到的结果。

\section{实验清单与论文落点}
本节给出可直接执行的实验清单。
每一项包含实验内容、实施方法、实验目的与论文呈现位置。

\subsection{A 必做实验\quad 主结果必须完整}

\subsubsection{A0 管线验证（Sanity Check，baseline\_debug）}

\paragraph{实验内容}
在进入任何正式对比（baseline\_full）之前，先完成一轮最小闭环验证（baseline\_debug）。
验证范围覆盖数据读取、训练与测试执行、指标计算、可视化导出、以及可复现快照。
当前已完成 RD4AD@RESC 的 baseline\_debug，后续每个数据集与每个基线方法均需先过本步骤。

\paragraph{成功判据}
\begin{itemize}
  \item 训练流程正常结束并生成权重文件。
  \item 测试流程正常结束并导出至少一份核心指标文件，例如图像级 AUROC。
  \item 导出至少一组异常热力图或定位可视化示例。
  \item 保存可复现快照，包含运行指令、依赖版本、代码提交记录与关键日志。
\end{itemize}

\paragraph{实施方法}
\begin{itemize}
  \item debug 训练轮数设为 1 到 5，仅验证闭环可用。
  \item 固定随机种子与运行环境标识，记录 code commit。
  \item 产物目录必须包含 checkpoints、metrics、figures、logs。
  \item 快照文件建议包含 env\_python.txt、env\_torch.txt、requirements\_freeze.txt、git\_commit.txt、git\_status.txt。
\end{itemize}

\paragraph{论文呈现}
baseline\_debug 不进入主结果表。
只作为复现实验链路证明，放附录或补充材料一段文字即可。

\subsubsection{A1 BMAD 六数据集主表（baseline\_full）}

\paragraph{实验内容}
在 BMAD 六个数据集上完成完整对比，形成一张主结果表。
像素级数据集为 RESC、BraTS、Liver。
图像级数据集为 OCT2017、RSNA、Camelyon16。

\paragraph{实施方法}
\begin{itemize}
  \item 每个数据集先完成 A0 的 baseline\_debug 验收，再进入 baseline\_full。
  \item baseline\_full 阶段锁定全部训练与评估设置，避免中途改动导致不可比。
  \item 每个数据集统一输入尺寸与预处理流程，全部使用 BMAD 官方 pipeline。
  \item 每个方法统一训练轮数与评估脚本，保证比较公平。
  \item 每个设置运行 3 个随机种子，报告 mean 与 std，并保留逐次运行结果与日志。
  \item 指标口径统一：
  \begin{itemize}
    \item 图像级数据集：报告 Image-level AUROC。
    \item 像素级数据集：报告 Image-level AUROC、Pixel-level AUROC、PRO。
  \end{itemize}
\end{itemize}

\paragraph{实验目的}
建立跨域主结论，并把复现成本压到可控范围。
用统一脚本与统一口径，减少审稿人对公平性的质疑。

\paragraph{论文呈现}
六数据集主结果表放正文。
完整训练细节、额外指标与实现说明放附录。

\subsubsection{A2 强基线对齐（baseline\_full，至少 6 个）}

\paragraph{实验内容}
在每个数据集上，除 RD4AD 外补齐足够强的 BMAD 内置基线。
建议覆盖三类方法，每类不少于 2 个，总数不少于 6 个。
所有基线都必须先完成 A0 的 baseline\_debug 验收，再纳入 baseline\_full。

\paragraph{基线建议}
\begin{itemize}
  \item 特征检索类：包含 PaDiM 与 PatchCore 的同类方法。
  \item 蒸馏或重建类：在 RD4AD 之外增加至少 1 个同类强基线。
  \item 一类学习或正则化类：从 BMAD 已集成方法中选择 1 到 2 个代表方法。
\end{itemize}

\paragraph{实施方法}
\begin{itemize}
  \item 对每个基线，先跑 baseline\_debug，确认 checkpoints、metrics、figures、logs 与快照齐全。
  \item 对进入主表的 baseline\_full，锁定输入分辨率、训练轮数、评测脚本与随机种子集合。
  \item run 命名必须包含 dataset、model、variant、seed、host，便于批量统计与追溯。
\end{itemize}

\paragraph{实验目的}
避免只与弱基线对比，避免只覆盖单一路线。
提升主结论可信度，同时降低复现与排错成本。

\paragraph{论文呈现}
与 A1 合并为同一张主表，或在主表中加入 Baselines 分组说明。
完整命令与配置快照放附录。

\subsection{B 主实验\quad 引入粒球机制（GBR 主方法）}
本节保留 B0 的内容不变，并在其后将原有 B1 以后内容升级为 \gbrvtwo 的正式主线。

\subsubsection{B0 方法定义与实现边界（GBR-AD，BMAD 插件化）}

\paragraph{实验内容}
实现一个粒球表示学习异常检测方法，记为 \textbf{GBR-AD}。
该方法包含四个模块。
粒球表示与图构建模块、异常检测模块、定位与分割读出模块、证据与可解释性模块。

\paragraph{方法要点}
\begin{itemize}
  \item \textbf{粒球生成（从粗到细）}：对输入做轻度平滑，计算梯度，优先从低梯度位置启动区域生长。
  通过纯度与方差阈值约束控制最细粒度。
  \item \textbf{粒球节点表征}：每个粒球用其覆盖区域的统计量描述。
  例如中心、尺度与均值方差等。
  可用如下统计定义中心与半径：
  \begin{equation}
    C=\frac{1}{n}\sum_{i=1}^{n}x_i,\qquad
    r=\frac{1}{n}\sum_{i=1}^{n}\lVert x_i-C\rVert .
  \end{equation}
  \item \textbf{图构建}：以“重叠即连边”为主规则。
  可选加入邻接边与包含边，形成多关系图。
  \item \textbf{异常打分与反投影}：训练阶段只用正常样本拟合正常图分布。
  测试阶段输出图像级异常分数，以及节点级分数。
  节点级分数按覆盖关系反投影回像素域生成 anomaly map。
  重叠区域采用加权融合，得到连续热图。
\end{itemize}

\paragraph{两种输入路径（必须都做）}
\begin{itemize}
  \item \textbf{GRI 路径}：直接在原图上生成粒球图，再做后续学习。
  \item \textbf{GRFM 路径}：先提取预训练 CNN 特征图，再在特征图上生成粒球图并学习。
  该路径通常收敛更快，泛化更强，作为默认主配置。
\end{itemize}

\paragraph{成功判据}
\begin{itemize}
  \item GBR-AD 在 RESC 上跑通 train 和 test，产物齐全。
  \item 粒球数量统计稳定，能输出每张图的节点数分布与平均重叠度。
  \item 能导出三类可视化：粒球划分叠加图、证据子图、反投影热图。
\end{itemize}

\paragraph{论文呈现}
方法结构图与伪代码放方法章节。
粒球划分与证据热图示例放可视化小节。

\subsubsection{B1 版本边界与主线调整（冻结 A 与 B0，升级 B1 以后内容）}

\paragraph{版本边界}
从本版 guideline 开始，GBR-AD 实验按两个版本管理。
\begin{itemize}
  \item \textbf{GBR-AD v1 provisional}：指已经完成的 B0 到当前 B1-v1 结果。
  这些结果用于证明插件化接入、可训练闭环和初步机制信号，不直接进入最终主表。
  \item \textbf{GBR-AD v2 official}：指本 guideline 新定义的方法与实验主线。
  从本版之后生成的 run、表格、图像与结论，均以 v2 为准。
\end{itemize}

\paragraph{冻结规则}
\begin{itemize}
  \item A 组结果与 B0 结果全部保留，不再改写。
  \item 已完成的 GBR-AD v1 结果保留为 provisional，不删除，不混入最终主表。
  \item 从新的 B1 开始，所有 proposed 结果统一按 v2 方法、统一 stopping policy、统一训练预算与统一评测脚本重跑。
\end{itemize}

\paragraph{论文呈现}
正文中的正式 proposed 结果一律使用 v2。
若需要说明工程演化过程，可在附录增加一段“从 B0 stub 到 B1 v2 的实现演化说明”。

\subsubsection{B2 GBR-AD v2 方法定义、完整数理逻辑与架构（正式主方法）}

\paragraph{方法总述}
\gbrvtwo 的核心思想是：以粒球而非像素或固定网格 token 作为主表示单元，以粒球图而非单一卷积感受野作为主结构单元，以粒球节点异常分数而非后验像素热图作为主检测单元。
CNN 或 ViT 特征提取器只提供稳定表征底座，图推理负责在粒球之间传播结构证据，图像级检测、像素级定位与证据链解释均从粒球图读出。

\paragraph{模块划分}
\gbrvtwo 按论文口径划分为五个模块。
\begin{enumerate}
  \item 多尺度粒球生成模块。
  \item 异构粒球图构建模块。
  \item 粒球图推理与正常性建模模块。
  \item 双读出模块，即图像级读出与像素级读出。
  \item 证据子图与可解释性模块。
\end{enumerate}

\paragraph{输入路径与统一主链}
给定输入图像 $x\in\mathbb{R}^{H\times W\times C}$，定义两条输入路径。
\begin{itemize}
  \item \gri：
  \begin{equation}
    x \rightarrow \mathcal{B}^{raw} \rightarrow G \rightarrow \{a_i\} \rightarrow \big(S_{img},A(u),G_E\big).
  \end{equation}
  \item \grfm：
  \begin{equation}
    x \xrightarrow{\phi_\theta} \{F^{(1)},F^{(2)},F^{(3)}\}
    \rightarrow \mathcal{B}^{feat}
    \rightarrow G
    \rightarrow \{a_i\}
    \rightarrow \big(S_{img},A(u),G_E\big).
  \end{equation}
\end{itemize}
其中，$\mathcal{B}^{raw}$ 与 $\mathcal{B}^{feat}$ 分别表示原图粒球集合与特征图粒球集合，$G$ 表示粒球异构图，$a_i$ 表示粒球节点异常分数，$S_{img}$ 表示图像级异常分数，$A(u)$ 表示像素级 anomaly map，$G_E$ 表示证据子图。

\paragraph{B2.1 多尺度粒球生成}
对每个尺度 $s\in\{1,\dots,S\}$，在原图或特征图上生成粒球候选区域 $\Omega_i^{(s)}$。
先对输入做轻度平滑，再计算梯度，并从低梯度区域开始区域生长。
记第 $s$ 个尺度上的第 $i$ 个粒球为
\begin{equation}
  g_i^{(s)}=\big(\Omega_i^{(s)},c_i^{(s)},r_i^{(s)},\mu_i^{(s)},\Sigma_i^{(s)},b_i^{(s)}\big),
\end{equation}
其中 $c_i^{(s)}$ 是中心，$r_i^{(s)}$ 是半径，$\mu_i^{(s)}$ 与 $\Sigma_i^{(s)}$ 分别是区域统计，$b_i^{(s)}$ 表示边界感知统计。
中心与半径定义为
\begin{equation}
  c_i^{(s)}=\frac{1}{|\Omega_i^{(s)}|}\sum_{u\in\Omega_i^{(s)}}p_u,
  \qquad
  r_i^{(s)}=\frac{1}{|\Omega_i^{(s)}|}\sum_{u\in\Omega_i^{(s)}}\lVert p_u-c_i^{(s)}\rVert_2.
\end{equation}

为保留“粒计算是灵魂”的方法定位，v2 显式引入粒球质量函数与分裂准则。
定义粒球质量为
\begin{equation}
  Q\big(g_i^{(s)}\big)=
  \alpha\,\mathrm{Compact}\big(g_i^{(s)}\big)
  +\beta\,\mathrm{EdgeAware}\big(g_i^{(s)}\big)
  +\gamma\,\mathrm{ScalePrior}\big(g_i^{(s)}\big),
\end{equation}
并对每个候选粒球施加如下递归细化规则：
\begin{equation}
  \text{split } g_i^{(s)}
  \quad \text{if}\quad
  \mathcal{P}\big(g_i^{(s)}\big)<\tau_p
  \;\text{or}\;
  \mathrm{tr}\big(\Sigma_i^{(s)}\big)>\tau_v,
\end{equation}
其中 $\tau_p$ 为纯度阈值，$\tau_v$ 为方差阈值。
粗粒球负责压缩背景与汇聚全局上下文，细粒球负责边界、小病灶与低对比局部结构，形成真正的多粒度粒球金字塔。

\paragraph{B2.2 节点表示与多尺度编码}
每个粒球节点的初始表示定义为
\begin{equation}
  h_i^{(0)}=[\mu_i,\sigma_i,c_i,r_i,s_i,b_i],
\end{equation}
其中 $\mu_i$ 与 $\sigma_i$ 为覆盖区域的统计特征，$c_i$ 与 $r_i$ 表示几何中心与半径，$s_i$ 表示尺度编码，$b_i$ 表示边界感知统计。
在 \gri 中，$\mu_i$ 与 $\sigma_i$ 来自原图区域统计；在 \grfm 中，$\mu_i$ 与 $\sigma_i$ 来自多层 CNN 特征图的区域统计。随后通过节点编码器映射到图推理空间：
\begin{equation}
  z_i^{(0)}=\psi\big(h_i^{(0)}\big),
\end{equation}
其中 $\psi(\cdot)$ 为轻量节点编码器，例如两层 MLP 或线性投影加归一化。

\paragraph{B2.3 异构粒球图构建}
定义粒球异构图为
\begin{equation}
  G=(V,E,R),
\end{equation}
其中 $V=\{g_i^{(s)}\}$ 是多尺度粒球节点集合，$R$ 为边关系集合。
v2 至少包含四类边：
\begin{enumerate}
  \item overlap edge：粒球覆盖区域有重叠即连边。
  \item adjacency edge：粒球边界距离低于阈值 $\delta$ 即连边。
  \item containment edge：细粒球与其上层父粒球之间建立包含边。
  \item semantic-kNN edge：在节点表示空间内按 kNN 建立语义近邻边。
\end{enumerate}
四类边对应四类结构作用：连续区域一致性、边界传播、粗细尺度协同、跨区域语义聚合。

\paragraph{B2.4 关系感知图推理}
图推理层定义为 relation-aware 多尺度消息传递：
\begin{equation}
  z_i^{(\ell+1)}
  =z_i^{(\ell)}+
  \sum_{r\in R}
  \psi_r\Bigg(
  \sum_{j\in\mathcal{N}_r(i)}
  \alpha_{ij}^{(r)}W_r z_j^{(\ell)}
  \Bigg),
\end{equation}
其中 $R$ 为边类型集合，$\mathcal{N}_r(i)$ 表示关系 $r$ 下与节点 $i$ 相邻的节点集合，$\alpha_{ij}^{(r)}$ 是关系感知注意力，$W_r$ 是不同关系的线性变换。
该层的目标不是单纯做图分类，而是在粒球之间传播结构证据，使粗粒球获得全局上下文，细粒球保留边界与小病灶线索，并让跨尺度关系显式进入异常检测逻辑。

\paragraph{B2.5 正常性建模与双分数机制}
训练阶段仅使用正常样本。
设图推理后的节点表示为 $z_i$，并定义每个尺度上的正常原型库或正常分布模型为 $\mathcal{M}_{normal}^{(s)}$。
v2 的节点异常分数不再只依赖单一重建误差，而定义为双分数与结构残差的组合：
\begin{equation}
  a_i=
  \lambda_d\,d\big(z_i,\mathcal{M}_{normal}^{(s)}\big)
  +\lambda_r\,\lVert \hat{h}_i-h_i\rVert_1
  +\lambda_c\,\Delta_i^{ctx}
  +\lambda_e\,\Delta_i^{edge}.
\end{equation}
其中：
\begin{itemize}
  \item $d(z_i,\mathcal{M}_{normal}^{(s)})$ 表示节点与正常原型库或正常高斯分布之间的距离。
  \item $\lVert \hat{h}_i-h_i\rVert_1$ 表示解码器重建误差。
  \item $\Delta_i^{ctx}$ 表示图上下文不一致性残差，用于抑制孤立噪声。
  \item $\Delta_i^{edge}$ 表示边关系违约项，用于刻画结构不连续或跨尺度不一致。
\end{itemize}
该设计的作用是：原型距离更偏向 image-level detection，重建误差更偏向局部异常响应，上下文残差与边关系残差帮助降低误报并增强结构一致性。

\paragraph{B2.6 训练目标}
对正常训练样本，定义总损失为
\begin{equation}
  \mathcal{L}=\eta_r\,\mathcal{L}_{rec}
  +\eta_p\,\mathcal{L}_{proto}
  +\eta_c\,\mathcal{L}_{ctx}
  +\eta_e\,\mathcal{L}_{edge},
\end{equation}
其中：
\begin{align}
  \mathcal{L}_{rec}&=\frac{1}{|V|}\sum_i\lVert \hat{h}_i-h_i\rVert_1,\\
  \mathcal{L}_{proto}&=\frac{1}{|V|}\sum_i d\big(z_i,\mathcal{M}_{normal}^{(s)}\big),\\
  \mathcal{L}_{ctx}&=\frac{1}{|V|}\sum_i\Delta_i^{ctx},\\
  \mathcal{L}_{edge}&=\frac{1}{|E|}\sum_{(i,j,r)\in E}\Delta_{ij}^{(r)}.
\end{align}
训练阶段的目标是让正常样本在多尺度粒球图上形成稳定、紧凑且结构一致的正常表示，从而使异常样本在测试时表现为节点级偏离。

\paragraph{B2.7 图像级读出}
图像级异常分数定义为三部分证据聚合：
\begin{equation}
  S_{img}=\beta_1\,\mathrm{TopKMean}(a)
  +\beta_2\,\mathrm{SubgraphMass}(G_E)
  +\beta_3\,\mathrm{ScaleConsensus}(a).
\end{equation}
其中：
\begin{itemize}
  \item $\mathrm{TopKMean}(a)$ 表示最高异常节点的 top-k 平均分数。
  \item $\mathrm{SubgraphMass}(G_E)$ 表示证据子图整体异常强度。
  \item $\mathrm{ScaleConsensus}(a)$ 表示异常证据在多尺度上的一致性。
\end{itemize}
该读出强调“结构证据链”而非单一高分节点，因此更有利于稳定的图像级检测。

\paragraph{B2.8 像素级反投影与边界感知读出}
节点分数按覆盖关系反投影到像素域。
对任意像素 $u$，其 anomaly score 定义为
\begin{equation}
  A(u)=
  \frac{\sum_{i:u\in\Omega_i}w_{u,i}a_i}
       {\sum_{i:u\in\Omega_i}w_{u,i}},
\end{equation}
其中权重定义为
\begin{equation}
  w_{u,i}=w_{u,i}^{dist}\cdot w_{u,i}^{edge}\cdot w_i^{purity}\cdot w_i^{scale}.
\end{equation}
各项含义如下：
\begin{itemize}
  \item $w^{dist}$：像素到粒球中心越近，权重越大。
  \item $w^{edge}$：边界区域更信任细粒球与边界感知节点。
  \item $w^{purity}$：纯度高、方差受控的粒球更可信。
  \item $w^{scale}$：小半径粒球在像素级定位中权重更大，避免大粒球稀释小病灶与边界。
\end{itemize}
为进一步提升边界恢复质量，可在最终热图后接一个轻量细化头，但该细化头只允许作为粒球反投影的局部边界校正器，不能替代粒球主表示与主读出链条。

\paragraph{B2.9 证据子图与可解释性输出}
定义高异常节点集合为
\begin{equation}
  V_E=\{i\mid a_i>\tau_a\},
\end{equation}
对应证据子图为
\begin{equation}
  G_E=G[V_E].
\end{equation}
最终系统输出四类结果：
\begin{itemize}
  \item 图像级分数 $S_{img}$。
  \item 像素级 anomaly map $A(u)$。
  \item 证据子图 $G_E$。
  \item 粒球划分叠加图、证据热图与失败模式可视化。
\end{itemize}
该模块保证方法不仅给出“是否异常”，还给出“哪些粒球支撑了结论，以及这些粒球如何在图上相互关联”。

\paragraph{B2.10 v2 的方法学定位}
\gbrvtwo 的核心创新不在于“粒球后面加一个普通 GNN”，而在于：
\begin{enumerate}
  \item 用多尺度粒球取代单尺度像素或固定 patch 作为主表示单元。
  \item 用异构粒球图显式编码结构关系，而不是只依赖卷积感受野隐式建模。
  \item 用双分数机制将全局 detection 与局部 localization 解耦并重新耦合。
  \item 用边界感知反投影把粒球级异常证据稳定地恢复到像素域。
  \item 用证据子图把性能、机制与可解释性联结成统一证据链。
\end{enumerate}

\paragraph{论文呈现}
本小节对应方法章节主体。
建议在正文中配套给出一张总架构图、一张多尺度粒球图构建图、一张证据子图示意图，以及一段算法伪代码。

\subsubsection{B3 主对比实验（GBR-AD v2 正式主线）}

\paragraph{实验内容}
B1 以后所有正式 proposed 结果均按 \gbrvtwo 执行。
主对比对象包含：RD4AD、GBR-AD(GRFM-v2)、GBR-AD(GRI-v2)。
正式实验不再混用 v1 与 v2 结果。

\paragraph{阶段划分}
\begin{enumerate}
  \item \textbf{B3.1 RESC 方法验证阶段}：仅在 RESC 上比较 RD4AD、GRFM-v2 与 GRI-v2，目的是快速验证 v2 是否同时改善 image-level 与 pixel-level。
  \item \textbf{B3.2 像素级三数据集阶段}：先完成 RESC、BraTS、Liver 的三方法对比，主攻定位能力与 PRO。
  \item \textbf{B3.3 六数据集正式阶段}：在像素级阶段结论稳定后，再扩展到 OCT2017、RSNA、Camelyon16，形成最终六数据集主表。
  \item \textbf{B3.4 最终冻结阶段}：锁定最终配置后，不再改动训练策略、后处理与评测脚本，统一生成正文主表与附录表。
\end{enumerate}

\paragraph{统一训练设置建议}
\begin{itemize}
  \item 训练只用正常样本。
  \item 统一使用 3 个随机种子，推荐 $\{0,1,42\}$。
  \item 统一采用 plateau early stop，max\_epochs 设为 100，上限固定。
  \item segmentation 任务以 Pixel-level AUROC 作为 early stop 监控量，image-level 任务以 Image-level AUROC 作为监控量。
  \item 统一优化器、学习率与权重衰减，不允许在 baseline 与 proposed 之间使用不同 stopping policy、不同测试脚本或不同阈值策略。
\end{itemize}

\paragraph{主表指标}
\begin{itemize}
  \item 像素级数据集：Image-level AUROC、Pixel-level AUROC、PRO。
  \item 图像级数据集：Image-level AUROC。
  \item 所有 GBR-AD 额外报告：平均粒球数、平均边数、推理耗时、显存峰值。
\end{itemize}

\paragraph{补充指标建议}
为提升论文层次，v2 附录建议补充以下指标，但不替代主表口径。
\begin{itemize}
  \item Image-level AP。
  \item Image-level F1-max。
  \item Pixel-level AP。
  \item Pixel-level F1-max。
  \item AU-PRO0.05 或低误报预算下的 PRO 变体。
  \item Dice 与 mIoU，仅限阈值选择严格基于验证集时使用。
\end{itemize}

\paragraph{实验目的}
形成 \gbrvtwo 的正式主结论，并回答以下核心问题：
\begin{enumerate}
  \item 与 RD4AD 相比，粒球机制是否在至少部分 BMAD 数据集上取得更高检测或定位指标。
  \item 在 \grfm 与 \gri 两条输入路径中，哪一条更适合作为主配置。
  \item 粒球带来的图结构与多粒度优势，是否值得其额外计算代价。
\end{enumerate}

\paragraph{论文呈现}
正文包含六数据集主表、Proposed 主实验表与效率表。
附录补充完整 3-seed 表格、补充指标与逐数据集分析。

\subsubsection{B4 受控识别实验（困难子集专测，v2 强化版）}

\paragraph{实验内容}
在像素级三数据集上构造四类困难子集。
\begin{enumerate}
  \item \textbf{小病灶子集}：按 GT mask 面积占比排序，取最小一档。
  \item \textbf{低对比子集}：按病灶与邻域背景的标准化对比度排序，取最低一档。
  \item \textbf{细边界子集}：按 boundary complexity 或 perimeter-to-area ratio 排序，取最高一档。
  \item \textbf{低误报预算子集}：固定较低像素误报预算后比较 lesion recall、PRO 或 fixed-FPR TPR。
\end{enumerate}

\paragraph{实施方法}
\begin{itemize}
  \item 子集划分只使用 GT 与输入统计，不改动训练数据与训练流程。
  \item 全量结果与困难子集结果共用同一模型输出，保证比较公平。
  \item 对每个子集同时报告主指标与固定误报预算下的召回或 PRO 变体。
\end{itemize}

\paragraph{实验目的}
直接验证 H4，即粒球机制是否对边界细节、小病灶与低对比区域具有更强表达力。
使“为什么 v2 有效”可以从困难场景而非平均场景中被识别出来。

\paragraph{论文呈现}
正文至少放 1 张困难子集对比表或 1 张多曲线图。
附录给出划分阈值与子集统计量。

\subsubsection{B5 证据链可解释性（正向证据与失败模式并列）}

\paragraph{实验内容}
对 RESC、BraTS、Liver 的代表样本，输出以下两类内容。
\begin{enumerate}
  \item \textbf{正向证据链}：输入图像、粒球划分叠加、证据子图、证据热图、预测掩膜五联图。
  \item \textbf{失败模式分类}：误检与漏检案例，按机制原因分类，例如大粒球稀释、小病灶未被细化、边界过平滑、图传播过度平滑、反投影权重失衡等。
\end{enumerate}

\paragraph{实施方法}
\begin{itemize}
  \item 用节点分数与关系注意力选取 Top-K 节点构建证据子图。
  \item 统计证据子图覆盖 GT 的比例、证据子图平均节点数与边数。
  \item 对失败案例给出机制标签，而不是只给图像示意。
\end{itemize}

\paragraph{实验目的}
强化“可审查证据”而非“只有分数”的方法学风格。
让方法的性能、机制与可解释性形成闭环。

\paragraph{论文呈现}
正文放 1 张综合大图。
附录放更多案例和失败模式图库。

\subsection{C 必做实验\quad 证明贡献来自核心模块}

\subsubsection{C1 关键消融实验（v2 版，至少 6 组）}

\paragraph{实验内容}
在 RESC 与 BraTS 上优先完成以下消融。
\begin{enumerate}
  \item 去掉多尺度粒球，只保留单尺度粒球。
  \item 去掉异构多关系边，只保留 overlap 边。
  \item 去掉原型距离分支，只保留重建误差。
  \item 去掉边界感知反投影，只保留简单平均反投影。
  \item 去掉轻量细化头，只保留粒球原始反投影。
  \item 去掉证据子图质量项，只保留 top-k 节点平均读出。
\end{enumerate}

\paragraph{实施方法}
\begin{itemize}
  \item 每次只改一个因素，其余训练设置保持一致。
  \item 每组消融运行 3 个随机种子并报告 mean 与 std。
  \item 主消融优先在 RESC 与 BraTS 上完成，后续再补 Liver。
\end{itemize}

\paragraph{实验目的}
证明 v2 的提升来自粒球与图结构本身，而不是训练技巧、阈值碰巧或 backbone 偶然更强。

\paragraph{论文呈现}
消融主表放正文。
更多补充消融放附录。

\subsubsection{C2 超参敏感性（聚焦 v2 核心超参）}

\paragraph{实验内容}
仅改变 1 到 2 个关键超参，输出性能曲线或小表。
优先覆盖：
\begin{itemize}
  \item 粒球纯度阈值 $\tau_p$。
  \item 粒球方差阈值 $\tau_v$。
  \item 节点数上限 $N_{max}$。
  \item 小粒球反投影权重系数。
\end{itemize}

\paragraph{实验目的}
展示方法对关键粒球超参不脆弱，并给出默认建议范围。

\paragraph{论文呈现}
放入附录一页。
正文用一两句话概述并引用附录图表。

\subsection{D 强烈建议实验\quad 医学场景加分项}

\subsubsection{D1 数据效率}

\paragraph{实验内容}
在 RESC、BraTS、Liver 上做训练数据比例实验。
训练集只使用正常样本，分别取 10\%、25\%、50\%、100\%。
其余设置不变。

\paragraph{实验目的}
突出粒球先验在低数据条件下的稳定性与泛化优势。

\paragraph{论文呈现}
给出一张曲线图或小表，放正文补充实验或附录。

\subsubsection{D2 鲁棒性测试}

\paragraph{实验内容}
对测试集施加轻度扰动，不修改标注。
扰动包括噪声、模糊、亮度变化、对比度变化、小角度旋转与轻微缩放。

\paragraph{实验目的}
验证方法在噪声、设备差异和轻度域偏移下的稳定性。

\paragraph{论文呈现}
给出一张鲁棒性曲线图或小表。

\subsubsection{D3 跨分辨率与跨设备稳健性}

\paragraph{实验内容}
保持训练配置尽量不变，只改变输入分辨率或设备来源。
至少在 1 到 2 个代表数据集上做跨分辨率或跨设备测试。

\paragraph{实验目的}
验证粒球机制是否天然具有尺度适应性与成像条件稳健性。

\paragraph{论文呈现}
正文一张小表或附录一张对比表。

\subsection{E 必做实验\quad 代价、可解释性与归档}

\subsubsection{E1 计算开销与结构统计}

\paragraph{实验内容}
报告参数量与推理耗时，至少给出每张图推理时间。
同时报告显存峰值或可用 batch size 上限。
对 GBR-AD v2 额外报告以下结构统计：
\begin{itemize}
  \item 每张图平均粒球数。
  \item 每张图平均边数。
  \item 细粒球占比。
  \item 证据子图平均节点数。
  \item 证据子图平均边数。
\end{itemize}

\paragraph{实验目的}
若方法引入粒球和图结构，需要说明其代价是否合理，并解释结构复杂度与性能变化之间的对应关系。

\paragraph{论文呈现}
效率表可放正文或附录。
正文建议用一行概述性能与效率的权衡。

\subsubsection{E2 定性可视化与错误分析}

\paragraph{实验内容}
在 RESC、BraTS、Liver 每个数据集选择 8 到 12 张典型样本。
展示输入图像、GT mask、预测 anomaly map、证据子图与二值化结果。
同时选择 4 到 6 张失败案例，按误检与漏检分类并分析原因。

\paragraph{实验目的}
医学异常检测强调定位可解释性。
定性结果与错误分析提升可信度与可审查性。

\paragraph{论文呈现}
定性图放正文。
更多案例放附录。

\subsubsection{E3 归档规范}

\paragraph{每个 run 必须保存的内容}
\begin{itemize}
  \item checkpoints/
  \item metrics/
  \item figures/
  \item logs/
  \item env\_python.txt
  \item env\_torch.txt
  \item requirements\_freeze.txt
  \item git\_commit.txt
  \item git\_status.txt
  \item run\_command.txt
\end{itemize}

\paragraph{命名规范}
所有正式对比 run 使用统一命名模板：
\begin{equation}
  \texttt{YYYYMMDD-HHMMSS\_\_P-BMAD\_\_T-anom\_\_D-<dataset>\_\_M-<model>\_\_V-<variant>\_\_S-<seed>\_\_H-<host>}
\end{equation}

\section{推荐推进顺序（Revised after B0）}
建议按如下顺序推进，以较小成本形成完整证据。
\begin{enumerate}
  \item 保留 A 组与 B0 已完成结果，不再改写。
  \item 将当前已经得到的 GBR-AD v1 结果统一标记为 \textbf{v1 provisional}。
  \item 先在 RESC 上实现并验证 \gbrvtwo，包括 GRFM-v2 与 GRI-v2 两条路径。
  \item 通过 RESC 的方法验证后，先完成像素级三数据集，即 RESC、BraTS、Liver 的主对比。
  \item 补齐图像级三数据集，即 OCT2017、RSNA、Camelyon16，形成六数据集主表。
  \item 在主表稳定后，执行 B4 的困难子集实验与 B5 的证据链可解释性实验。
  \item 完成 C 组消融，优先验证多尺度粒球、异构边、双分数与边界感知反投影的贡献。
  \item 补齐 D 与 E，形成鲁棒性、泛化性、效率、证据链与归档闭环。
\end{enumerate}

\section{论文交付物清单（v2 正式版）}
最终需要交付的图表与材料如下。
\begin{itemize}
  \item 主结果表 1 张，覆盖六数据集与全部方法，报告 mean 与 std。
  \item Proposed 主实验表 1 张，至少包含 RD4AD、GBR-AD(GRFM-v2)、GBR-AD(GRI-v2)。
  \item 困难子集对比表或曲线 1 张，覆盖小病灶、低对比、细边界与低误报预算专测。
  \item 消融表 1 张，至少 6 行关键消融设置。
  \item 鲁棒性图或小表 1 张。
  \item 数据效率图或小表 1 张。
  \item 效率表 1 张，包含推理耗时、显存、粒球数、边数、细粒球占比与证据子图统计。
  \item 定性可视化图 1 到 2 张，包含粒球划分、证据子图、证据热图与失败案例分析。
  \item 方法架构图 1 张，明确显示多尺度粒球生成、异构图构建、双分数检测与边界感知反投影。
  \item 附录中的完整训练设置表、补充指标表、归档说明与 run-to-table 映射表。
\end{itemize}

\end{document}